\documentclass[12pt,a4paper]{article}
\usepackage[utf8]{inputenc}
\usepackage[T1]{fontenc}
\usepackage{amsmath}
\usepackage{amsfonts}
\usepackage{amssymb}
\usepackage{graphicx}
\usepackage{hyperref}
\usepackage{listings}
\usepackage{xcolor}
\usepackage{geometry}
\geometry{margin=1in}

\title{Infinity Mode Protocol Extension\\RFC 25.02201\\Game Event Notification Protocol Extension}
\author{Tootsville Development Team}
\date{\today}

\begin{document}

\maketitle

\begin{abstract}
This RFC proposes an extension to the Infinity Mode protocol to formalize the handling of game event notifications. The extension establishes standardized message formats and distribution mechanisms for in-game events, ensuring proper separation between game events and system-level notifications. This addresses the need to replace ad-hoc messaging systems with a protocol-compliant approach while maintaining backward compatibility with existing client implementations.
\end{abstract}

\section{Introduction}

The current Infinity Mode protocol provides robust mechanisms for real-time communication between clients and servers in the Tootsville game environment. However, the protocol lacks specific provisions for handling game event notifications in a standardized manner. This RFC proposes an extension that:

\begin{itemize}
\item Defines standardized message formats for game events
\item Establishes clear separation between game events and system notifications
\item Provides mechanisms for event routing and prioritization
\item Ensures protocol compliance for all game-related communications
\end{itemize}

\section{Historical Context}

\subsection{ActionScript Implementation}

The original ActionScript 3.0 implementation utilized several notification mechanisms:

\begin{itemize}
\item \texttt{Tootsville.Gossip.Parrot.say()} - Used for various game events and system messages
\item \texttt{Event} system - Flash Player's native event dispatching
\item Custom event classes like \texttt{WorldEvent} for specific game occurrences
\item Direct manipulation of UI components for immediate feedback
\end{itemize}

This approach, while functional, suffered from:
\begin{itemize}
\item Lack of clear separation between game events and system notifications
\item Inconsistent message formatting
\item Tight coupling between game logic and UI presentation
\item Difficulty in maintaining event routing and prioritization
\end{itemize}

\subsection{Java Server Implementation}

The Java server (\texttt{AppiusClaudiusCaecus}) handled events through:

\begin{itemize}
\item Command processing system with reflection-based dispatch
\item Datagram-based communication for real-time events
\item Threaded event handling with metronome-based game loops
\item Centralized logging and error reporting
\end{itemize}

\section{Protocol Extension Specification}

\subsection{Message Format}

All game event notifications SHALL use the following JSON message format:

\begin{lstlisting}[language=JSON, caption=Game Event Message Format]
{
  "type": "game-event",
  "event-id": "unique-event-identifier",
  "timestamp": "ISO-8601-timestamp",
  "priority": "normal|high|low",
  "category": "equipment|player|system|npc|world",
  "target": {
    "type": "user|all|area",
    "user-id": "optional-user-identifier",
    "area-id": "optional-area-identifier"
  },
  "event": {
    "type": "effect|status|interaction|notification",
    "subtype": "specific-event-type",
    "data": {
      // Event-specific payload
    }
  },
  "metadata": {
    "source": "equipment|player|server|world",
    "persistent": false,
    "ttl": 30000
  }
}
\end{lstlisting}

\subsection{Event Categories}

\subsubsection{Equipment Events}

Equipment-related events include item usage, effects, and status changes:

\begin{itemize}
\item \texttt{equipment-effect} - Direct effects from equipment usage
\item \texttt{equipment-status} - Status updates (cooldown, empty, active)
\item \texttt{equipment-interaction} - Interactions between equipment items
\end{itemize}

\subsubsection{Player Events}

Player-centric events involve character actions and state changes:

\begin{itemize}
\item \texttt{player-action} - Player-initiated actions
\item \texttt{player-status} - Player state changes
\item \texttt{player-interaction} - Interactions between players
\end{itemize}

\subsubsection{System Events}

System-level events for game state management:

\begin{itemize}
\item \texttt{system-notification} - General system announcements
\item \texttt{system-maintenance} - Maintenance-related notifications
\item \texttt{system-error} - Error conditions requiring user attention
\end{itemize}

\subsubsection{NPC Events}

Events involving non-player characters:

\begin{itemize}
\item \texttt{npc-interaction} - NPC-initiated interactions
\item \texttt{npc-behavior} - NPC behavior changes
\item \texttt{npc-spawn} - NPC creation/destruction events
\end{itemize}

\subsubsection{World Events}

World-level events affecting the game environment:

\begin{itemize}
\item \texttt{world-change} - Environmental changes
\item \texttt{world-event} - Scheduled world events
\item \texttt{world-transition} - Area transitions and loading
\end{itemize}

\subsection{Message Routing}

\subsubsection{Client-Side Routing}

Clients SHALL route messages based on the following priority order:

\begin{enumerate}
\item High-priority events (immediate user attention required)
\item Normal-priority events (standard game feedback)
\item Low-priority events (background information)
\end{enumerate}

\subsubsection{Server-Side Distribution}

Servers SHALL distribute messages using the following mechanisms:

\begin{itemize}
\item \textbf{WebSocket}: Real-time delivery for active connections
\item \textbf{HTTP Polling}: Fallback for clients without WebSocket support
\item \textbf{Push Notifications}: Mobile and desktop notifications
\item \textbf{Email/SMS}: Critical system notifications (limited use)
\end{itemize}

\subsection{Event Persistence}

\subsubsection{Transient Events}

Most game events are transient and SHALL NOT be persisted beyond immediate delivery:

\begin{itemize}
\item Equipment effect notifications
\item Player interaction feedback
\item NPC behavior updates
\end{itemize}

\subsubsection{Persistent Events}

Certain events MAY require persistence for later retrieval:

\begin{itemize}
\item Achievement notifications
\item Quest progress updates
\item System maintenance announcements
\end{itemize}

\subsection{Error Handling}

\subsubsection{Message Delivery Failures}

In case of delivery failures:

\begin{enumerate}
\item Attempt redelivery with exponential backoff
\item Queue messages for later delivery if client reconnects
\item Discard messages after TTL expiration
\item Log delivery failures for debugging
\end{enumerate}

\subsubsection{Invalid Message Format}

Clients and servers SHALL:

\begin{itemize}
\item Validate message format against schema
\item Log validation errors
\item Discard malformed messages
\item Continue processing valid messages
\end{itemize}

\section{Implementation Guidelines}

\subsection{Client Implementation}

\subsubsection{Vue/TypeScript Implementation}

The Vue/TypeScript client SHALL implement the following components:

\begin{lstlisting}[language=TypeScript, caption=GameEventManager Interface]
// GameEventManager.ts
interface GameEventMessage {
  type: 'game-event';
  eventId: string;
  timestamp: string;
  priority: 'normal' | 'high' | 'low';
  category: 'equipment' | 'player' | 'system' | 'npc' | 'world';
  target: {
    type: 'user' | 'all' | 'area';
    userId?: string;
    areaId?: string;
  };
  event: {
    type: 'effect' | 'status' | 'interaction' | 'notification';
    subtype: string;
    data: any;
  };
  metadata: {
    source: string;
    persistent: boolean;
    ttl: number;
  };
}

class GameEventManager {
  private eventQueue: GameEventMessage[] = [];
  private subscribers: Map<string, Function[]> = new Map();

  public send(message: GameEventMessage): void {
    // Implementation for sending messages
  }

  public subscribe(category: string, handler: Function): void {
    // Implementation for subscribing to events
  }

  public processQueue(): void {
    // Implementation for processing queued messages
  }
}
\end{lstlisting}

\subsubsection{Message Processing}

The client SHALL process messages in the following order:

\begin{enumerate}
\item Validate message format
\item Check user permissions for message visibility
\item Route message to appropriate handlers
\item Update UI state if necessary
\item Log message for debugging
\end{enumerate}

\subsection{Server Implementation}

\subsubsection{Common Lisp Implementation}

The server SHALL implement message routing using the following structure:

\begin{lstlisting}[language=Lisp, caption=Server Event Router]
(defclass game-event-router ()
  ((event-queue :initform (make-instance 'queue))
   (subscribers :initform (make-hash-table :test 'equal))
   (message-validator :initform (make-instance 'message-validator))))

(defmethod route-event ((router game-event-router) event)
  "Route a game event to appropriate subscribers"
  (when (validate-event (message-validator router) event)
    (let ((category (getf event :category)))
      (dolist (subscriber (gethash category (subscribers router)))
        (send-to-subscriber subscriber event)))))

(defmethod broadcast-event ((router game-event-router) event)
  "Broadcast event to all connected clients in relevant area"
  (dolist (client (find-clients-in-area event))
    (websocket-send client (encode-event event))))
\end{lstlisting}

\section{Backward Compatibility}

\subsection{Transition Period}

During the transition period:

\begin{itemize}
\item Legacy \texttt{Tootsville.Gossip.Parrot.say()} calls SHALL be deprecated
\item New game event system SHALL be implemented alongside legacy system
\item Clients SHALL support both old and new message formats
\item Gradual migration of existing code to new system
\end{itemize}

\subsection{Legacy Message Mapping}

Legacy messages SHALL be mapped to new format as follows:

\begin{lstlisting}[language=JSON, caption=Legacy Message Mapping]
{
  "legacy": "Tootsville.Gossip.Parrot.say('Player pushed!', 'push')",
  "new-format": {
    "type": "game-event",
    "category": "equipment",
    "event": {
      "type": "effect",
      "subtype": "push",
      "data": { "message": "Player pushed!" }
    }
  }
}
\end{lstlisting}

\section{Security Considerations}

\subsection{Message Validation}

All messages SHALL be validated to prevent:

\begin{itemize}
\item Message injection attacks
\item Unauthorized event broadcasting
\item Client-side message spoofing
\item Resource exhaustion through message flooding
\end{itemize}

\subsection{Access Control}

Event visibility SHALL be controlled based on:

\begin{itemize}
\item User authentication status
\item User permissions and roles
\item Geographic/game area restrictions
\item Event sensitivity level
\end{itemize}

\section{Performance Considerations}

\subsection{Message Throughput}

The system SHALL support:

\begin{itemize}
\item 1000+ messages per second during peak usage
\item Sub-100ms latency for high-priority events
\item Efficient message queuing and processing
\item Memory-efficient message storage
\end{itemize}

\subsection{Client Resource Usage}

Clients SHALL:

\begin{itemize}
\item Limit concurrent message processing
\item Implement message rate limiting
\item Provide mechanisms for message batching
\item Handle large message volumes gracefully
\end{itemize}

\section{Testing and Validation}

\subsection{Unit Testing}

Comprehensive unit tests SHALL cover:

\begin{itemize}
\item Message format validation
\item Event routing logic
\item Error handling scenarios
\item Performance benchmarks
\end{itemize}

\subsection{Integration Testing}

Integration tests SHALL verify:

\begin{itemize}
\item End-to-end message delivery
\item Multi-client synchronization
\item Server failover scenarios
\item Network interruption handling
\end{itemize}

\subsection{Load Testing}

Load tests SHALL validate:

\begin{itemize}
\item Maximum concurrent users
\item Peak message throughput
\item Memory usage under load
\item Database performance
\end{itemize}

\section{Deployment Strategy}

\subsection{Phased Rollout}

The extension SHALL be deployed in phases:

\begin{enumerate}
\item \textbf{Phase 1}: Server-side implementation and testing
\item \textbf{Phase 2}: Client-side implementation with legacy support
\item \textbf{Phase 3}: Migration of existing equipment and game logic
\item \textbf{Phase 4}: Legacy system deprecation and removal
\end{enumerate}

\subsection{Rollback Plan}

In case of deployment issues:

\begin{itemize}
\item Maintain legacy system as fallback
\item Implement feature flags for gradual rollout
\item Monitor error rates and performance metrics
\item Prepare rollback scripts and procedures
\end{itemize}

\section{Future Extensions}

\subsection{Potential Enhancements}

Future versions MAY include:

\begin{itemize}
\item Message encryption for sensitive events
\item Advanced filtering and search capabilities
\item Event replay and debugging features
\item Integration with external notification services
\end{itemize}

\section{References}

\begin{enumerate}
\item Infinity Mode Protocol Specification (Base Protocol)
\item ActionScript 3.0 Implementation Documentation
\item Java Server Architecture Documentation
\item Vue/TypeScript Client Implementation Guide
\item WebSocket Protocol Specification (RFC 6455)
\end{enumerate}

\end{document}